\documentclass[11pt]{article}

\usepackage[margin=1in]{geometry}
\usepackage{amsmath, amssymb, amsthm}
\usepackage{booktabs}
\usepackage{microtype}
\usepackage{natbib}
\usepackage[hidelinks]{hyperref}

% --- notation ---
\newcommand{\R}{\mathbb{R}}
\newcommand{\Sig}{\Sigma}
\newcommand{\F}{\mathcal{F}}
\newcommand{\B}{\mathcal{B}}
\newcommand{\Aset}{\mathcal{A}}
\newcommand{\Sset}{\mathcal{S}}
\newcommand{\trans}{^{\top}}
\newcommand{\lb}{\ell}
\newcommand{\ub}{u}
\newcommand{\inv}{^{-1}}

\newtheorem{proposition}{Proposition}

\title{The Critical Line Algorithm \emph{is} a constrained LASSO:\\
a dictionary between two homotopies}

\author{%
  [Authors]\thanks{Draft for \emph{The American Statistician}. The companion
  software paper describes \texttt{cvxcla}, the open implementation used for the
  numerical checks below.}%
}

\date{Draft, \today}

\begin{document}
\maketitle

\begin{abstract}
Markowitz's Critical Line Algorithm (CLA), which traces the mean--variance
efficient frontier, and the LARS/LASSO homotopy, which traces the
$\ell_1$-regularised regression path, are taught in different fields and seldom
presented together. They are the same algorithm. Both follow the optimiser of a
strictly convex quadratic as a single scalar parameter varies, both produce a
piecewise-linear path, and both trace it by an active-set walk from one breakpoint
to the next. We give the exact dictionary between the two, and then push past the
textbook (unconstrained) correspondence: a \emph{single} active-set homotopy traces
the Markowitz frontier and the LASSO path under \emph{arbitrary} linear equality and
inequality constraints, with box bounds, non-negativity, and a gross-exposure
($\ell_1$) cap as special cases. The unifying lens is that one side carves its
active set with explicit polyhedral constraints while the other has it induced by
the $\ell_1$ subgradient; the gross-exposure portfolio is exactly where the two
coincide. One implementation runs both, and reproduces \texttt{scikit-learn}'s LASSO
path to machine precision.
\end{abstract}

\section{Two paths, one mechanism}

A surprising number of methods compute a \emph{path}: not a single fit, but the
whole family of optimisers as one scalar parameter is swept. Two of the most used
live in different literatures. In finance, the Critical Line Algorithm
\citep{markowitz1956, markowitz1959} traces the mean--variance efficient frontier as
the risk-aversion parameter varies. In statistics, least angle regression and the
LASSO homotopy \citep{tibshirani1996, osborne2000, efron2004} trace the
$\ell_1$-regularised regression coefficients as the penalty varies. They are taught
apart, cited apart, and implemented apart.

Our point is that they are one algorithm. This is not a loose analogy: the two
problems are instances of the same parametric quadratic program, their solution
maps are the same piecewise-linear object, and the active-set walk that traces one
traces the other verbatim. Making the correspondence precise does three things. It
lets a reader fluent in one field read the other for free; it transfers
machinery (an exact, constrained LASSO path falls out of the portfolio side); and
it isolates what is genuinely particular to each problem from what is shared
scaffolding.

\section{The two parametric quadratic programs}

\paragraph{The portfolio frontier.}
Let $w\in\R^n$ be portfolio weights, $\mu\in\R^n$ expected returns, and
$\Sig\in\R^{n\times n}$ a symmetric positive (semi)definite covariance. The CLA
traces the \emph{return-parametrised} quadratic program
\begin{equation}
\label{eq:cla}
\min_{w\in\R^n}\ \tfrac12\, w\trans \Sig\, w \;-\; \lambda\, \mu\trans w
\quad\text{subject to}\quad \mathbf{1}\trans w = 1,\ \ \lb \le w \le \ub,
\end{equation}
as $\lambda$ decreases from $+\infty$ (the maximum-return portfolio) to $0$ (the
global minimum-variance portfolio). The locus of optimisers is the efficient
frontier.

\paragraph{The regression path.}
For a response $y\in\R^m$ and design $X\in\R^{m\times n}$, the LASSO solves
\begin{equation}
\label{eq:lasso}
\min_{\beta\in\R^n}\ \tfrac12\,\|y - X\beta\|_2^2 \;+\; \lambda\,\|\beta\|_1 ,
\end{equation}
and its minimiser $\beta(\lambda)$ is traced as $\lambda$ decreases from the value
$\lambda_{\max}=\|X\trans y\|_\infty$ (where $\beta=0$) towards $0$.

\paragraph{Why both paths are piecewise linear.}
Fix the \emph{active set}. For \eqref{eq:cla} this is the \emph{free} set $\F$ of
assets strictly inside their bounds (the rest pinned at $\lb$ or $\ub$); for
\eqref{eq:lasso} it is the support $\Aset=\{j:\beta_j\neq 0\}$ together with the
signs $s_\Aset=\operatorname{sign}(\beta_\Aset)$. On the region of $\lambda$ where
the active set is constant, the stationarity (KKT) conditions are linear in the
unknowns with a \emph{constant} system matrix and a right-hand side
\emph{affine} in $\lambda$:
\begin{align}
\label{eq:cla-seg}
\text{(CLA)}\quad & \Sig_{\F\F}\, w_\F + \nu\,\mathbf{1}_\F = \lambda\,\mu_\F - \Sig_{\F\B}w_\B,
   && \mathbf{1}\trans w = 1, \\
\label{eq:lasso-seg}
\text{(LASSO)}\quad & X_\Aset\trans X_\Aset\, \beta_\Aset = X_\Aset\trans y - \lambda\, s_\Aset .
\end{align}
Hence each solves to an affine function of $\lambda$,
\[
w_\F(\lambda) = \alpha + \lambda\,\beta, \qquad
\beta_\Aset(\lambda) = a - \lambda\, c ,
\]
valid until the active set must change. Those change-points are the CLA's
\emph{turning points} and the LASSO's \emph{breakpoints}. This is the
\citet{rosset2007} criterion at work: a quadratic loss together with a polyhedral
(piecewise-linear) penalty or constraint forces a piecewise-linear path. In the
language of parametric optimisation both are \emph{one-parameter parametric
quadratic programs}, whose solution map is piecewise affine over a polyhedral
partition of $\lambda$; the same structure underlies explicit model predictive
control \citep{bemporad2002}.

\section{The dictionary}

Comparing \eqref{eq:cla-seg} and \eqref{eq:lasso-seg}, the covariance $\Sig$ plays
the role of the Gram matrix $X\trans X$ and the expected-return vector $\mu$ the
role of $X\trans y$. Under that single substitution the two homotopies line up term
for term (Table~\ref{tab:dict}). The walk is identical: at each step compute the
affine segment, find the smallest decrease in $\lambda$ at which an active-set event
occurs, apply it, repeat. Even the anti-cycling rule is shared: both inherit the
simplex method's lowest-index (Bland) tie-break to stay deterministic on degenerate
problems.

\begin{table}[t]
\centering
\begin{tabular}{lll}
\toprule
Ingredient & Critical Line Algorithm & LARS / LASSO homotopy \\
\midrule
Swept parameter & risk aversion $\lambda:\infty\!\to\!0$ & penalty $\lambda:\infty\!\to\!0$ \\
Quadratic form  & covariance $\Sig$ & Gram matrix $X\trans X$ \\
Linear term     & expected returns $\mu$ & correlations $X\trans y$ \\
Solution path   & $w(\lambda)=\alpha+\lambda\beta$ & $\beta(\lambda)=a-\lambda c$ \\
Active set      & free assets (off the bounds) & support (nonzero $\beta_j$) \\
Enter event     & blocked multiplier changes sign & correlation reaches $\lambda$ \\
Leave event     & free weight reaches a box bound & coefficient reaches zero \\
Step length     & \multicolumn{2}{l}{smallest positive step to the next breakpoint} \\
Degeneracy      & lowest-index tie-break & lowest-index tie-break \\
\bottomrule
\end{tabular}
\caption{The Critical Line Algorithm and the LARS/LASSO homotopy as two instances of
one parametric active-set scheme. The only substitution is
$(\Sig,\mu)\leftrightarrow(X\trans X,\,X\trans y)$.}
\label{tab:dict}
\end{table}

There is one genuinely instructive difference, and it is structural rather than
cosmetic. In the CLA the active set is carved out by \emph{explicit} polyhedral
constraints: an asset is blocked when it sits on a box bound, and the events are
weights reaching bounds or multipliers releasing them. In the LASSO the active set
is \emph{induced} by the $\ell_1$ penalty through its subgradient: sparsity is an
emergent property of the geometry, not a constraint the modeller writes down. The
same polyhedral mechanism runs in opposite directions, one imposing the faces and
the other discovering them. (A second, minor difference: forward LARS only adds
variables, whereas the CLA, like the full LASSO homotopy, both adds and drops them
as $\lambda$ decreases.)

\section{Beyond the textbook case: general linear constraints}

The correspondence above is the unconstrained (or budget-only) one. Its real
payoff is that it survives \emph{general} linear constraints, and that the
extension is the same on both sides. Augment either problem with an equality system
$A x = b$ and inequalities $G x \le h$:
\begin{equation}
\label{eq:general}
\min_x\ \tfrac12\, x\trans H x - \lambda\, d\trans x \ \ (\text{or}\ +\lambda\|x\|_1)
\quad\text{s.t.}\quad A x = b,\ \ G x \le h,\ \ \lb \le x \le \ub,
\end{equation}
with $(H,d)=(\Sig,\mu)$ for the frontier and $(H,d)=(X\trans X, X\trans y)$ for the
path. An \emph{active} inequality row, held at $g_i\trans x = h_i$, enters the
reduced KKT system exactly as an equality row does: stacking
$C=[\,A;\,G_{\Sset}\,]$ over the active rows $\Sset$ recovers the equality-only
solve with $A$ replaced by $C$. So each step is still one solve against the
free/support block of $H$, bordered by an $(m+|\Sset|)\times(m+|\Sset|)$ Schur
complement. Box bounds and coefficient signs stay a separate, cheaper per-variable
active set rather than being folded into $G$, so the common constraint-free problem
pays nothing for the generality. (It is worth recording the natural wrong turn:
turning an inequality into an equality with a non-negative slack leaves the slack
with no curvature, so the augmented Hessian is only positive \emph{semi}definite and
the reduced solve is singular whenever the inequality is inactive. The bordered
system avoids this.)

This matters because the statistics side rarely has an exact \emph{constrained}
path. Coordinate descent and proximal methods deliver an approximate path on a
$\lambda$-grid; LARS gives the unconstrained path. Equation~\eqref{eq:general}
traces, exactly and in one pass, the LASSO path under arbitrary $G\beta \le h$, under
non-negativity $\beta\ge 0$ (the analogue of a long-only portfolio), and under a
linear equality. The dictionary now reads: long-only $\leftrightarrow$ non-negative
LASSO; the budget $\mathbf{1}\trans w = 1$ $\leftrightarrow$ a linear equality on
$\beta$; a sector- or factor-neutrality row $\leftrightarrow$ a linear equality;
a group-exposure cap $\leftrightarrow$ an inequality row.

\paragraph{Where the two literally coincide.}
The bound-versus-penalty duality is not only conceptual. \citet{brodie2009} add a
gross-exposure constraint $\|w\|_1 \le c$ to the Markowitz problem and show the
result \emph{is} a LASSO, with $c$ promoting sparse, stable portfolios. For a
long-only, fully invested portfolio the constraint is silent, since the budget
already fixes $\|w\|_1 = \mathbf{1}\trans w = 1$; this is why the plain CLA traces a
frontier rather than a sparsity path. Once short positions are admitted under a
gross-exposure cap, the explicit constraint and the induced penalty describe the
same polyhedron, and the parametric program the CLA traces and the LASSO path are
not merely similar but identical.

\section{One implementation, two problems}

Because the walk is shared, a single path-tracing engine serves both: it touches
the quadratic form $H$ only through a matrix--vector product, a solve against the
active block, and an active-to-inactive cross product, and is told only what an
``event'' means. The \texttt{cvxcla} package
\citep{cvxcla} implements exactly this. Its \texttt{Lasso} object feeds the Gram
matrix $X\trans X$ and the vector $X\trans y$ to the same loop that traces the
portfolio frontier; the protocol it calls is named for a generic quadratic form,
not for the covariance, for this reason.

Two checks, deliberately light, confirm the identification.

\emph{The unconstrained path is the LASSO path.} On a standard regression design,
\texttt{cvxcla} and \texttt{scikit-learn}'s \texttt{lars\_path} report the same
number of breakpoints and the same coefficients to machine precision: across the
design we use, the maximum coefficient discrepancy is
$\max_j |\hat\beta_j^{\texttt{cvxcla}} - \hat\beta_j^{\texttt{sklearn}}| = 3.2\times
10^{-15}$ over $13$ breakpoints. The two implementations trace the identical path;
they differ only in heritage.

\emph{The constrained path is exact.} Under a per-coordinate non-negativity
constraint $\beta\ge 0$ and a linear inequality $G\beta\le h$, the path traced by
\eqref{eq:general} matches a reference quadratic-program solver
(\texttt{CVXPY}/Clarabel) solved independently at each breakpoint to a maximum
coefficient discrepancy of $\approx 6\times 10^{-10}$ -- solver precision -- whereas
a grid-based approximation only samples the same curve.

\section{What is, and is not, new}

The ingredients are individually known, and we are explicit about which.
\citet{efron2004} established the LARS/LASSO homotopy; \citet{osborne2000} gave an
earlier homotopy view. \citet{rosset2007} characterised exactly when a regularised
path is piecewise linear, a result that subsumes both problems. \citet{brodie2009}
identified the gross-exposure Markowitz problem as a LASSO. The parametric-QP and
explicit-MPC viewpoint \citep{bemporad2002} is classical in control. What we add is
the consolidation that these pieces invite but do not state: a single active-set
homotopy, and an explicit term-for-term dictionary, under which the
\emph{constrained} Markowitz frontier and the \emph{constrained} LASSO path (general
linear equalities and inequalities, non-negativity, and the gross-exposure ball) are
the same computation -- extending the unconstrained correspondence and the single
gross-exposure special case to arbitrary polyhedral constraints -- together with the
reference implementation that traces both.

\section{Discussion}

Reading the two algorithms as one is more than tidy. It tells a portfolio
practitioner that the machinery for sparse, leverage-constrained allocations is the
LASSO homotopy, and it hands a statistician an exact path under constraints that the
usual solvers approximate or cannot express. It also clarifies pedagogy: the
critical line and the LASSO path are the same object seen from the constraint side
and the penalty side, and the gross-exposure portfolio is the hinge between them.
The broader moral is the one parametric optimisation has long known and the two
fields rediscovered separately: when a strictly convex quadratic is followed along a
single parameter under polyhedral structure, the path is piecewise linear and an
active set is all the state you need to trace it exactly.

\begin{thebibliography}{Markowitz(1959)}

\bibitem[Bemporad et~al.(2002)]{bemporad2002} Bemporad, A., Morari, M., Dua, V., and
Pistikopoulos, E.~N. (2002). The explicit linear quadratic regulator for constrained
systems. \emph{Automatica}, 38(1), 3--20.

\bibitem[Brodie et~al.(2009)]{brodie2009} Brodie, J., Daubechies, I., De Mol, C.,
Giannone, D., and Loris, I. (2009). Sparse and stable Markowitz portfolios.
\emph{Proceedings of the National Academy of Sciences}, 106(30), 12267--12272.

\bibitem[Authors(2025)]{cvxcla} Authors (2025). \texttt{cvxcla}: the Critical Line
Algorithm with a covariance-operator interface. Companion software paper.

\bibitem[Efron et~al.(2004)]{efron2004} Efron, B., Hastie, T., Johnstone, I., and
Tibshirani, R. (2004). Least angle regression. \emph{The Annals of Statistics},
32(2), 407--499.

\bibitem[Markowitz(1956)]{markowitz1956} Markowitz, H.~M. (1956). The optimization of
a quadratic function subject to linear constraints. \emph{Naval Research Logistics
Quarterly}, 3(1--2), 111--133.

\bibitem[Markowitz(1959)]{markowitz1959} Markowitz, H.~M. (1959). \emph{Portfolio
Selection: Efficient Diversification of Investments}. Wiley, New York.

\bibitem[Osborne et~al.(2000)]{osborne2000} Osborne, M.~R., Presnell, B., and
Turlach, B.~A. (2000). A new approach to variable selection in least squares
problems. \emph{IMA Journal of Numerical Analysis}, 20(3), 389--403.

\bibitem[Rosset and Zhu(2007)]{rosset2007} Rosset, S. and Zhu, J. (2007). Piecewise
linear regularized solution paths. \emph{The Annals of Statistics}, 35(3),
1012--1030.

\bibitem[Tibshirani(1996)]{tibshirani1996} Tibshirani, R. (1996). Regression
shrinkage and selection via the lasso. \emph{Journal of the Royal Statistical
Society, Series B}, 58(1), 267--288.

\end{thebibliography}

\end{document}
